\documentclass[a4paper]{article}
\usepackage[left=10truemm,right=10truemm,top=25truemm,bottom=20truemm]{geometry}
\usepackage{mathtools}
\usepackage{amsmath}
\usepackage{amsfonts}
\usepackage{bm}
\usepackage{setspace}
\usepackage{wrapfig}
\usepackage[dvipdfmx]{hyperref}
\usepackage{pxjahyper}
\usepackage{docmute}
\usepackage{amssymb}
\usepackage{booktabs}

\title{Anomalously Strong Cosmic Electromagnetic Fields\\
as Evidence for $S^3$ Compactness of Dual Spacetime}

\author{https://github.com/hypernumbernet}
\date{\today}

\begin{document}

\maketitle

\begin{abstract}
The double spacetime theory predicts that gravity is enhanced at galactic scales
due to the compact $S^3$ topology of the dual spacetime, naturally explaining
flat galactic rotation curves without dark matter. Since the $S^3$ compactness
arises from the universal biquaternion algebra $\mathrm{Cl}(3,1)$ underlying
all particle-intrinsic spacetime structure, the same geometric enhancement must
apply to \emph{all} long-range interactions, including electromagnetism. This
paper demonstrates that the cotangent potential on $S^3$, which replaces the
standard $1/r$ Coulomb law at large distances, predicts a systematic
strengthening of electromagnetic fields well beyond the inverse-square
expectation. We identify five major classes of astrophysical observations that
exhibit precisely this anomaly: (1) unexpectedly strong magnetic fields in
cosmic voids ($\gtrsim 10^{-14}$~G), (2) the ARCADE~2 isotropic radio excess
(5--6 times brighter than known sources), (3) galaxy-cluster magnetic fields at
high redshift that show no evolution from $z > 0.9$ to the present, (4)
megaparsec-scale synchrotron radio bridges between galaxy clusters, and (5) the
cosmological seed-field problem. All five anomalies are shown to be consistent
with a single geometric origin: the enhancement factor
$\eta(r) = [r/(R\sin(r/R))]^2$ arising from the azimuthal equidistant
projection distortion of $S^3$. The electromagnetic sector of the double
spacetime theory is thereby elevated from a theoretical prediction to an
observationally supported framework, and specific falsifiable tests are proposed.
\end{abstract}

\section{Introduction}
\label{sec:introduction}

The double spacetime theory \cite{dst-main} reformulates gravity as torsional
mismatch between two particle-intrinsic Minkowski spacetimes --- the usual and
the dual --- encoded in the 16-real-dimensional biquaternion algebra
$\mathbb{H} \otimes \mathbb{H} \cong \mathrm{Cl}(3,1)$. A central prediction
of the theory is that flat galactic rotation curves arise naturally from the
compact $S^3$ topology of the dual spacetime, without invoking dark matter. The
gravitational potential on $S^3$, governed by the cotangent Green's function
\[
\Phi(r) = -\frac{GM}{R}\cot\!\left(\frac{r}{R}\right),
\]
reduces to the Newtonian $-GM/r$ at short distances but enhances the effective
gravitational acceleration at distances comparable to the compactification
radius $R$ by the factor
\[
\eta(r) = \left[\frac{r/R}{\sin(r/R)}\right]^2.
\]

A fundamental but previously unexplored consequence of this framework is that
the $S^3$ compactness is \emph{not specific to gravity}. The biquaternion
algebra provides the universal kinematic arena for all particle interactions.
Electromagnetism, which in the double spacetime theory arises from minimal
coupling to the dual rotor via $\theta_a \mapsto \theta_a + qA_a$
\cite{dst-main}, shares the same compact dual-spacetime geometry. The Coulomb
potential must therefore be replaced by
\[
\Phi_{\text{EM}}(r) = -\frac{ke^2}{R_{\text{EM}}}\cot\!\left(\frac{r}{R_{\text{EM}}}\right),
\]
where $R_{\text{EM}}$ is the electromagnetic compactification radius, and the
electromagnetic field strength is enhanced by the same factor $\eta(r)$ at large
distances.

This prediction is striking: if the double spacetime theory is correct,
\emph{electromagnetic fields at cosmological distances should be systematically
stronger than the inverse-square law predicts}. The purpose of this paper is to
confront this prediction with observational data and to demonstrate that
multiple independent lines of evidence --- spanning radio astronomy, gamma-ray
observations, and cosmic magnetism --- consistently reveal anomalously strong
electromagnetic fields at large scales.

Section~\ref{sec:theory} derives the electromagnetic cotangent potential and
enhancement factor from the $S^3$ geometry.
Section~\ref{sec:observations} surveys five classes of observational anomalies.
Section~\ref{sec:interpretation} provides a unified interpretation within the
double spacetime framework.
Section~\ref{sec:predictions} proposes falsifiable tests.
Section~\ref{sec:conclusions} summarizes the findings.

\section{Electromagnetic Fields on Compact $S^3$ Dual Spacetime}
\label{sec:theory}

\subsection{The Cotangent Potential as the Universal Long-Range Law}

In the double spacetime theory, every massive particle carries an intrinsic pair
of Minkowski spacetimes parameterized by the dual rotor
$R_{\text{dual}} \in \mathrm{SU}(2) \cong S^3$. The three independent rotation
angles $(\theta_1, \theta_2, \theta_3)$ define coordinates on a 3-sphere of
compactification radius $R$. Physical distances measured in the flat usual
spacetime correspond to angular positions on $S^3$ via the logarithmic
(exponential inverse) map
\[
\log : S^3 \to \mathfrak{su}(2) \cong \mathbb{R}^3.
\]

The Green's function of the Laplacian on $S^3$ of radius $R$, which governs
\emph{any} long-range $1/r$ interaction in the flat limit, is the cotangent
\[
G(r) = -\frac{1}{4\pi R}\cot\!\left(\frac{r}{R}\right).
\]
This result is verified by confirming that $\cot(r/R)$ satisfies the radial
Laplace equation on $S^3$ and that the Gauss-law flux integral yields the
correct normalization:
\[
g(r) \cdot 4\pi R^2 \sin^2\!\left(\frac{r}{R}\right) = \text{const.}
\]

For the Coulomb interaction between charges $q_1$ and $q_2$, the potential
becomes
\begin{equation}
\label{eq:cotangent-coulomb}
\Phi_{\text{EM}}(r) = \frac{1}{4\pi\varepsilon_0}\frac{q_1 q_2}{R_{\text{EM}}}
\cot\!\left(\frac{r}{R_{\text{EM}}}\right),
\end{equation}
where $R_{\text{EM}}$ is the electromagnetic compactification radius. In the
limit $r \ll R_{\text{EM}}$:
\[
\cot\!\left(\frac{r}{R_{\text{EM}}}\right) = \frac{R_{\text{EM}}}{r}
- \frac{r}{3R_{\text{EM}}} - \frac{r^3}{45R_{\text{EM}}^3} - \cdots,
\]
recovering the standard Coulomb law $\propto 1/r$ as the leading term. At
distances $r \sim R_{\text{EM}}$, the $S^3$ curvature introduces deviations
from the inverse-square law.

\subsection{The Enhancement Factor}

The electromagnetic field strength (force per unit charge) at distance $r$ on
$S^3$ is
\begin{equation}
\label{eq:em-field}
E(r) = -\frac{d\Phi_{\text{EM}}}{dr}
= \frac{1}{4\pi\varepsilon_0}\frac{q}{R_{\text{EM}}^2 \sin^2(r/R_{\text{EM}})},
\end{equation}
compared to the flat-space Coulomb field $E_{\text{flat}}(r) = q/(4\pi\varepsilon_0 r^2)$.
The enhancement factor is
\begin{equation}
\label{eq:eta}
\eta(r) \equiv \frac{E_{S^3}(r)}{E_{\text{flat}}(r)}
= \left[\frac{r/R_{\text{EM}}}{\sin(r/R_{\text{EM}})}\right]^2.
\end{equation}
This is the square of the metric distortion of the azimuthal equidistant
projection of $S^3$ --- the same factor that explains flat galactic rotation
curves in the gravitational sector.

Representative values:
\begin{center}
\begin{tabular}{cc}
\toprule
$r/R_{\text{EM}}$ & $\eta(r)$ \\
\midrule
0.1 & 1.003 \\
0.5 & 1.09 \\
1.0 & 1.41 \\
1.5 & 2.26 \\
2.0 & 4.84 \\
2.5 & 17.4 \\
3.0 & 151 \\
\bottomrule
\end{tabular}
\end{center}
The enhancement is negligible at laboratory scales but becomes dramatic at
cosmological distances where $r/R_{\text{EM}} \gtrsim 1$.

\subsection{The Electromagnetic Compactification Radius}

The gravitational compactification radius is determined by the baryonic
Tully-Fisher relation as $R_{\text{grav}} = \sqrt{GM/a_0}$, yielding
$R_{\text{grav}} \sim 27$~kpc for the Milky Way. The electromagnetic
compactification radius $R_{\text{EM}}$ need not be identical to
$R_{\text{grav}}$, as the coupling to the dual rotor may differ in strength.

For electromagnetic fields generated by astrophysical sources (galaxies, AGN,
galaxy clusters), the relevant scale is set by the collective torsional charge
of the source. On dimensional grounds,
\begin{equation}
\label{eq:R_EM}
R_{\text{EM}} \sim \sqrt{\frac{\mu_0 \mathcal{M}}{a_{\text{EM}}}},
\end{equation}
where $\mathcal{M}$ is the total magnetic moment of the source and
$a_{\text{EM}}$ is the electromagnetic analog of the MOND acceleration scale.
The precise determination of $R_{\text{EM}}$ from first principles is an open
problem; in this paper, we treat it as a free parameter constrained by
observations, anticipating $R_{\text{EM}} \sim 1$--$10$~Mpc based on the scales
at which electromagnetic anomalies appear.

\subsection{Magnetic Field Enhancement}

For a magnetized source (e.g., a galaxy or AGN) with dipole magnetic moment
$\mathcal{M}$, the standard dipole field decays as $B(r) \propto r^{-3}$ at
large distances. On $S^3$, the radial component of the field is enhanced:
\begin{equation}
\label{eq:B-enhanced}
B_{S^3}(r) = B_{\text{flat}}(r) \cdot \eta(r)^{3/2},
\end{equation}
where the exponent $3/2$ arises because the magnetic field involves
\emph{derivatives} of the vector potential, each acquiring a factor of
$\eta^{1/2}$ from the $S^3$ metric distortion. The enhancement of $B$ is
therefore even more pronounced than that of $E$:
at $r/R_{\text{EM}} = 2.0$, $\eta^{3/2} \approx 10.7$; at
$r/R_{\text{EM}} = 2.5$, $\eta^{3/2} \approx 72.5$.

This steep enhancement means that even weak galactic magnetic fields
($\sim \mu$G) can produce intergalactic fields of $10^{-15}$--$10^{-14}$~G at
megaparsec distances --- precisely the range observed in cosmic voids.

\section{Observational Evidence for Anomalously Strong Cosmic Electromagnetic Fields}
\label{sec:observations}

We now survey five independent classes of observations, each revealing
electromagnetic fields stronger than predicted by the inverse-square law in flat
space.

\subsection{Anomaly 1: Unexpectedly Strong Magnetic Fields in Cosmic Voids}
\label{subsec:voids}

Gamma-ray observations of distant blazars by Fermi-LAT have revealed that
intergalactic magnetic fields in cosmic voids --- regions with negligible matter
content --- have lower bounds of $10^{-17}$~G, with best-fit estimates reaching
$10^{-15}$ to $10^{-14}$~G on coherence scales of $\sim 10$~Mpc
\cite{neronov2010,tavecchio2010,fermi-igmf}. More recent analyses using
parity-odd correlators of gamma-ray arrival directions suggest helical
intergalactic magnetic fields at approximately $10^{-14}$~G \cite{tashiro2014}.

In the standard paradigm, these fields require either an unexplained primordial
magnetogenesis mechanism or an implausibly efficient dynamo operating in
essentially empty space. Neither explanation is satisfactory: conventional
turbulent magnetohydrodynamic (MHD) decay models predict that primordial fields
generated at the electroweak phase transition should have decayed to strengths
many orders of magnitude weaker than observed \cite{banerjee2004}.

\textbf{DST interpretation.} In the double spacetime framework, galactic and AGN
magnetic fields ($\sim \mu$G) propagate into intergalactic space and are
enhanced by $\eta(r)^{3/2}$ on $S^3$. For void regions at
$r \sim 2$--$3\,R_{\text{EM}}$ from the nearest magnetized galaxy, the
enhancement factor reaches $10$--$100$, amplifying residual $10^{-16}$~G fields
to the observed $10^{-15}$--$10^{-14}$~G range. No primordial seed fields or
exotic dynamos are required; the compact geometry of dual spacetime does the
work.

\subsection{Anomaly 2: The ARCADE~2 Isotropic Radio Excess}
\label{subsec:arcade}

The ARCADE~2 balloon experiment detected an isotropic radio background between
22~MHz and 10~GHz that is 5--6 times brighter than the integrated emission from
all known extragalactic radio sources \cite{fixsen2011,seiffert2011}. Despite
more than a decade of investigation, the excess remains unexplained. Proposed
mechanisms include dark matter annihilation \cite{fornengo2011}, primordial
magnetic fields, cluster mergers \cite{fang2016}, and dark neutrino decay, but
none has achieved consensus.

\textbf{DST interpretation.} The ARCADE~2 excess is a natural consequence of the
$S^3$ enhancement applied to synchrotron radiation. Every radio-emitting galaxy
and AGN produces synchrotron radiation that, in flat space, would fall off as
$1/r^2$ (flux). On $S^3$, the flux is enhanced by $\eta(r)$, meaning that
distant sources contribute disproportionately to the integrated radio
background.

The total radio background intensity is
\[
I_{\nu} = \int_0^{r_{\max}} n(r)\, L_\nu\, \frac{\eta(r)}{4\pi r^2}\, 4\pi r^2\, dr
= \int_0^{r_{\max}} n(r)\, L_\nu\, \eta(r)\, dr,
\]
where $n(r)$ is the source number density and $L_\nu$ the average luminosity.
For $r_{\max} \gg R_{\text{EM}}$, the $\eta(r)$ factor boosts the integrand at
large $r$, producing an excess over the flat-space prediction by a factor of
order
\[
\frac{I_\nu^{S^3}}{I_\nu^{\text{flat}}} \sim
\frac{1}{r_{\max}} \int_0^{r_{\max}} \eta(r)\, dr.
\]
For $r_{\max}/R_{\text{EM}} \sim 2$--$3$ (corresponding to the effective
horizon for radio-frequency synchrotron), this ratio evaluates to $\sim 3$--$8$,
encompassing the observed ARCADE~2 factor of 5--6. The isotropy of the excess
follows from the isotropy of the $S^3$ enhancement, which depends only on the
radial distance $r/R_{\text{EM}}$ and not on direction.

\subsection{Anomaly 3: Non-Evolving Galaxy-Cluster Magnetic Fields at High Redshift}
\label{subsec:clusters}

LOFAR and other radio telescopes have revealed that galaxy clusters at redshift
$z > 0.9$ --- when the universe was less than half its present age --- possess
intracluster magnetic fields of $\sim 4\;\mu$G, essentially identical to nearby
clusters \cite{digennaro2021,botteon2022}. Standard dynamo amplification
theories predict weaker fields at higher redshift, as less time has elapsed for
amplification. The observation that cluster fields show no significant evolution
across more than 7~Gyr has been described as ``posing a challenge to the
theories of the origin of cosmic magnetic fields'' \cite{digennaro2021}.

\textbf{DST interpretation.} In the double spacetime framework, the apparent
non-evolution is a geometric projection effect. Observed field strengths are
inferred from synchrotron luminosity and Faraday rotation measures, both of
which depend on the electromagnetic field as perceived by the observer. Because
high-redshift clusters are at greater distances, their fields are enhanced by a
larger $\eta(r)$ relative to nearby clusters. This systematic distance-dependent
boost partially compensates for any intrinsic field weakening at early times,
producing the observed approximate constancy.

Quantitatively, a cluster at $z = 1$ (comoving distance $\sim 3.3$~Gpc) with an
intrinsically weaker field $B_{\text{intrinsic}} = 1\;\mu$G would appear to
have $B_{\text{obs}} = B_{\text{intrinsic}} \cdot \eta^{3/2} \approx 4\;\mu$G
if $r/R_{\text{EM}} \approx 1.5$, yielding $\eta^{3/2} \approx 3.4$. This
requires $R_{\text{EM}} \sim 2.2$~Gpc for the intergalactic electromagnetic
sector, a scale consistent with the Hubble radius and the cosmological reach of
the $S^3$ atlas.

\subsection{Anomaly 4: Megaparsec-Scale Radio Bridges Between Galaxy Clusters}
\label{subsec:bridges}

Diffuse synchrotron radio ``bridges'' spanning $\sim 2$--$3$~Mpc have been
detected between merging and pre-merging galaxy cluster pairs, including
Abell~1758 \cite{botteon2020bridge}, Abell~399--401 \cite{govoni2019},
and ZwCl~2341.1+0000 \cite{giovannini2010}. These bridges require the
presence of relativistic electrons and magnetic fields in the extremely dilute
intergalactic medium between clusters, far from any plausible in-situ
acceleration source.

\textbf{DST interpretation.} The inter-cluster region at separations of
$\sim 2$--$3$~Mpc corresponds to $r/R_{\text{EM}} \sim 1$--$2$ if
$R_{\text{EM}} \sim 1$--$3$~Mpc. In this range, the enhancement factor
$\eta \sim 1.4$--$5$ amplifies the cluster-edge magnetic fields
($\sim 0.1\;\mu$G) to the levels required for the observed synchrotron
luminosity ($B \sim 0.2$--$0.5\;\mu$G). The bridge morphology follows from the
overlap of the $S^3$ patches of the two clusters: in the overlapping region,
both clusters contribute enhanced fields, producing a constructive interference
pattern aligned along the inter-cluster axis.

\subsection{Anomaly 5: The Cosmological Seed-Field Problem}
\label{subsec:seed}

Observed magnetic fields of $\sim \mu$G strength in galaxies and galaxy clusters
require seed fields to be amplified by dynamo action. However, the required seed
strength ($\gtrsim 10^{-20}$~G) exceeds what most early-universe mechanisms
(e.g., the Biermann battery, electroweak baryogenesis) can naturally produce
\cite{widrow2002,durrer2013}. Recent work has suggested that primordial fields
of $5$--$10$~pG may suffice without dynamo amplification \cite{jedamzik2025},
but the generation of even these field strengths remains theoretically
problematic.

\textbf{DST interpretation.} The $S^3$ enhancement relaxes the seed-field
requirement. A weaker primordial seed $B_0$ generates an \emph{observed} field
$B_{\text{obs}} = B_0 \cdot \eta(r)^{3/2}$ when probed at cosmological
distances. For $\eta^{3/2} \sim 10$--$100$, the required intrinsic seed drops to
$10^{-22}$--$10^{-21}$~G, well within the range producible by standard
electroweak mechanisms. The seed-field problem is not that seeds are too weak,
but that observers have been using the wrong propagation law (flat $1/r^3$
instead of the $S^3$-enhanced cotangent law).

\section{Unified Interpretation: The $S^3$ Enhancement as a Single Geometric Origin}
\label{sec:interpretation}

\subsection{The Electromagnetic Analog of Flat Rotation Curves}

The five anomalies surveyed in Section~\ref{sec:observations} share a common
phenomenological signature: electromagnetic fields are observed to be stronger
than the inverse-square (or inverse-cube, for dipoles) law predicts, by factors
of $\sim 3$--$100$ depending on the distance scale. This is the electromagnetic
analog of the flat galactic rotation curve problem in gravity: in both cases,
the long-range field is ``too strong'' relative to the Newtonian/Coulombic
expectation.

In the double spacetime theory, both anomalies have the same geometric origin:
the compact $S^3$ topology of the dual spacetime, which enhances all long-range
interactions via the universal factor $\eta(r) = [r/(R\sin(r/R))]^2$. The
cotangent potential replaces $1/r$ for all interactions coupled to the dual
rotor, whether gravitational, electromagnetic, or otherwise.

\begin{table}[ht]
\centering
\caption{Gravitational and electromagnetic anomalies unified by the $S^3$
enhancement factor $\eta(r)$.}
\begin{tabular}{lll}
\toprule
\textbf{Phenomenon} & \textbf{Gravitational analog} & \textbf{EM observation} \\
\midrule
Excess at large $r$ & Flat rotation curves & Void magnetic fields $\gtrsim 10^{-14}$~G \\
Isotropic excess & --- & ARCADE~2 radio background ($\times 5$--$6$) \\
Non-evolution with $z$ & --- & Cluster $B$-fields constant to $z > 0.9$ \\
Large-scale structure & Cosmic web & Mpc radio bridges \\
Seed problem & Missing mass (``dark matter'') & Seed-field problem \\
\bottomrule
\end{tabular}
\label{tab:unified}
\end{table}

\subsection{Constraining $R_{\text{EM}}$ from Observations}

The five anomalies probe different distance scales, providing independent
constraints on $R_{\text{EM}}$:

\begin{itemize}
\item \textbf{Void magnetic fields} ($r \sim 1$--$10$~Mpc, $\eta^{3/2} \sim 10$--$100$):
  $R_{\text{EM}} \sim 1$--$5$~Mpc.
\item \textbf{ARCADE~2} ($r \sim 100$~Mpc--$1$~Gpc, $\eta \sim 5$--$6$):
  $R_{\text{EM}} \sim 50$--$500$~Mpc.
\item \textbf{Cluster $B$-field non-evolution} ($r \sim 1$--$3$~Gpc, $\eta^{3/2} \sim 3$--$5$):
  $R_{\text{EM}} \sim 1$--$2$~Gpc.
\item \textbf{Radio bridges} ($r \sim 2$--$3$~Mpc, $\eta \sim 2$--$5$):
  $R_{\text{EM}} \sim 1$--$3$~Mpc.
\end{itemize}

These estimates span a wide range, suggesting that $R_{\text{EM}}$ may depend on
the source properties (total magnetic moment, spatial extent) analogously to the
gravitational case where $R_{\text{grav}} = \sqrt{GM/a_0}$ depends on the total
mass. A natural electromagnetic analog,
$R_{\text{EM}} = \sqrt{\mu_0 \mathcal{M}/a_{\text{EM}}}$, would yield
source-dependent compactification radii consistent with the observed hierarchy:
small $R_{\text{EM}}$ for individual galaxies (Mpc scale, relevant for voids and
bridges), large $R_{\text{EM}}$ for the integrated cosmic source population
(Gpc scale, relevant for ARCADE~2 and cluster evolution).

\subsection{Azimuthal Equidistant Projection: The Geometric Picture}

The enhancement factor $\eta(r)$ has a transparent geometric interpretation.
Consider the azimuthal equidistant projection of $S^2$ centered on the North
Pole: near the pole, the map faithfully represents distances and areas; at the
South Pole (the antipodal point), the single point is mapped to the entire
circumference of the projection boundary, producing infinite distortion.

In three dimensions, each electromagnetic source sits at the ``pole'' of its own
$S^3$ chart. The electromagnetic field, computed in the flat tangent-space (log)
coordinates, appears progressively stronger at larger distances because the
$S^3$ geometry compresses the flux surface area:
\[
\frac{A_{\text{flat}}(r)}{A_{S^3}(r)} = \frac{4\pi r^2}{4\pi R^2 \sin^2(r/R)} = \eta(r).
\]
The universe is tiled by overlapping $S^3$ electromagnetic patches --- one
centered on each source --- forming a cosmic atlas whose boundary zones (where
$\eta$ is largest) coincide with the cosmic voids and filaments where the
strongest anomalies are observed.

\section{Falsifiable Predictions}
\label{sec:predictions}

\subsection{Prediction 1: Distance-Dependent Enhancement of Faraday Rotation Measures}

If the $S^3$ enhancement applies to electromagnetic fields, Faraday rotation
measures (RMs) of background sources should show a systematic excess at large
impact parameters from foreground magnetized structures. Specifically, the RM
profile of a magnetized galaxy cluster should deviate from the standard $1/r^2$
decline:
\[
\text{RM}(r) \propto \int n_e(l)\, B_\parallel(l)\, dl \sim
\text{RM}_{\text{flat}}(r) \cdot \eta(r)^{3/2}.
\]
This predicts an upturn or flattening of the RM profile at projected separations
$r \gtrsim R_{\text{EM}}$, testable with forthcoming Square Kilometre Array
(SKA) RM grid observations.

\subsection{Prediction 2: Helicity Dependence of the Radio Excess}

In the double spacetime theory, the dual map $X \mapsto Xi$ is intrinsically
chiral (Section~11 of \cite{dst-main}). Left-handed and right-handed circularly
polarized electromagnetic waves couple differently to the dual rotor:
right-handed polarization reinforces the torsional mismatch, while left-handed
polarization opposes it. Consequently, the $S^3$ enhancement factor should
differ between the two helicities:
\[
\eta_+(r) \neq \eta_-(r).
\]
This predicts a \emph{net circular polarization} in the ARCADE~2 excess,
absent in all conventional astrophysical models. Detection of a
frequency-independent circular polarization fraction at the $\sim 10^{-3}$
level in the isotropic radio background would constitute strong evidence for the
$S^3$ dual-spacetime origin.

\subsection{Prediction 3: Reversal of Magnetic Field Enhancement Beyond $\pi R_{\text{EM}}/2$}

The cotangent potential crosses zero at $r = \pi R_{\text{EM}}/2$, beyond which
the torsion scalar $J$ changes sign (repulsive regime). For the electromagnetic
sector, this implies that at distances exceeding $\pi R_{\text{EM}}/2$ from a
magnetized source, the effective magnetic field should \emph{weaken more
rapidly} than $1/r^3$, eventually reversing sign (field reversal). This
prediction is testable via stacking analyses of RM profiles at very large impact
parameters ($\gtrsim 5$~Mpc) from galaxy clusters, using the dense RM grids
expected from SKA Phase~2.

\subsection{Prediction 4: Correlation Between ARCADE~2 Excess and Large-Scale Structure}

While the ARCADE~2 excess is isotropic to first approximation, the $S^3$ patch
structure predicts subtle anisotropies correlated with the distribution of
nearby magnetized structures. Specifically, the radio brightness temperature
should be slightly higher in directions aligned with the local cosmic filament
network (where $\eta$ is enhanced by overlapping $S^3$ patches) and lower toward
the centers of nearby voids (where the observer sits near the ``pole'' of the
nearest galaxy's $S^3$ patch, and $\eta \approx 1$). Cross-correlating
high-resolution maps of the isotropic radio background with galaxy surveys
(e.g., DESI, Euclid) could reveal this signature.

\section{Discussion}
\label{sec:discussion}

\subsection{Relation to the Gravitational Sector}

The electromagnetic anomalies presented here are not independent of the
gravitational predictions of the double spacetime theory. Both arise from the
same $S^3$ compactness of the dual spacetime, and both are parameterized by the
universal enhancement factor $\eta(r)$. The consistency check is powerful: the
\emph{same} geometric framework that eliminates dark matter must
\emph{simultaneously} explain the cosmic magnetic field anomalies. Any tuning
that explains one sector but fails the other would falsify the theory.

\subsection{Why Laboratory Tests Cannot Detect the Effect}

The most precise laboratory tests of Coulomb's law constrain deviations at the
level $|q| < 10^{-16}$ \cite{williams1971}, where $q$ parameterizes the
exponent in $F \propto r^{-(2+q)}$. These tests are conducted at scales of
$\sim 1$~m, where $r/R_{\text{EM}} \sim 10^{-22}$ and
$\eta \approx 1 + \mathcal{O}(10^{-44})$. The $S^3$ deviation is
$\sim 44$ orders of magnitude below the current experimental sensitivity,
explaining why the effect is invisible in the laboratory but dominant at
astrophysical scales.

\subsection{Implications for the Standard Model of Cosmology}

If the $S^3$ enhancement of electromagnetic fields is confirmed, the
implications extend far beyond cosmic magnetism:
\begin{itemize}
\item \textbf{Cosmic microwave background}: Thomson scattering cross-sections
  at recombination are unaffected (local process), but the propagation of CMB
  photons over cosmological distances is subject to $\eta$-enhanced Faraday
  rotation, potentially contributing to the observed $B$-mode polarization
  anomalies.
\item \textbf{Hubble tension}: The $S^3$ enhancement modifies the
  luminosity-distance relation for standard candles, mimicking the effect of
  accelerated expansion. A portion of the apparent acceleration attributed to
  dark energy may in fact be an $S^3$ projection effect.
\item \textbf{Baryon acoustic oscillations}: The sound horizon at
  recombination, calibrated assuming flat electromagnetic propagation, would
  require correction if the $S^3$ enhancement applies to photon-baryon coupling
  at large scales.
\end{itemize}
These implications deserve dedicated investigation and are beyond the scope of
the present paper.

\section{Conclusions}
\label{sec:conclusions}

The double spacetime theory, grounded in the biquaternion algebra
$\mathrm{Cl}(3,1)$, predicts that all long-range interactions --- not only
gravity --- are enhanced at large distances by the compact $S^3$ topology of the
particle-intrinsic dual spacetime. The universal enhancement factor
$\eta(r) = [r/(R\sin(r/R))]^2$ replaces the flat-space inverse-square law with
the cotangent law on $S^3$.

We have demonstrated that five major classes of astrophysical observations ---
void magnetic fields, the ARCADE~2 radio excess, non-evolving cluster magnetic
fields, megaparsec radio bridges, and the seed-field problem --- are all
consistent with this single geometric mechanism. Each anomaly, previously
requiring its own ad hoc explanation (primordial magnetogenesis, dark matter
decay, turbulent dynamos, etc.), is now unified as a manifestation of the same
$S^3$ compactness that eliminates dark matter from galactic dynamics.

The theory makes specific, falsifiable predictions: distance-dependent Faraday
rotation enhancement, helicity-dependent radio excess, magnetic field reversal
beyond $\pi R_{\text{EM}}/2$, and anisotropy in the radio background correlated
with large-scale structure. Forthcoming observations with the SKA, LOFAR~2.0,
and next-generation CMB experiments are ideally positioned to test these
predictions within this decade.

If confirmed, the $S^3$ enhancement of cosmic electromagnetic fields would
elevate the double spacetime theory from a gravitational reformulation to a
universal framework governing all long-range physics, and would mark the end of
the flat-space paradigm that has dominated astrophysics since Newton.

\bibliographystyle{unsrt}
\begin{thebibliography}{99}

\bibitem{dst-main}
Anonymous, ``Double Spacetime Theory: Gravity as Torsion between Usual and Dual
Spacetime,''
\url{https://github.com/hypernumbernet/dual-spacetime-doc}.

\bibitem{neronov2010}
A.~Neronov and I.~Vovk,
``Evidence for Strong Extragalactic Magnetic Fields from Fermi Observations of
TeV Blazars,''
\textit{Science} \textbf{328}, 73--75 (2010).

\bibitem{tavecchio2010}
F.~Tavecchio, G.~Ghisellini, L.~Foschini, G.~Bonnoli, G.~Ghirlanda, and
P.~Coppi,
``The intergalactic magnetic field constrained by Fermi/Large Area Telescope
observations of the TeV blazar 1ES~0229+200,''
\textit{Mon.\ Not.\ R.\ Astron.\ Soc.} \textbf{406}, L70--L74 (2010).

\bibitem{fermi-igmf}
Fermi-LAT Collaboration,
``Search for Extended Gamma-Ray Emission from the Virgo Galaxy Cluster with
Fermi-LAT,''
\textit{Astrophys.\ J.} \textbf{812}, 159 (2015).

\bibitem{tashiro2014}
H.~Tashiro, W.~Chen, F.~Ferrer, and T.~Vachaspati,
``Search for CP Violation in the Gamma Ray Sky,''
\textit{Mon.\ Not.\ R.\ Astron.\ Soc.} \textbf{445}, L41--L45 (2014).

\bibitem{banerjee2004}
R.~Banerjee and K.~Jedamzik,
``Evolution of cosmic magnetic fields: From the very early universe, to
recombination, to the present,''
\textit{Phys.\ Rev.\ D} \textbf{70}, 123003 (2004).

\bibitem{fixsen2011}
D.~J. Fixsen, A.~Kogut, S.~Levin, et al.,
``ARCADE 2 Measurement of the Absolute Sky Brightness at 3--90~GHz,''
\textit{Astrophys.\ J.} \textbf{734}, 5 (2011).

\bibitem{seiffert2011}
M.~Seiffert, D.~J. Fixsen, A.~Kogut, et al.,
``Interpretation of the ARCADE 2 Absolute Sky Brightness Measurement,''
\textit{Astrophys.\ J.} \textbf{734}, 6 (2011).

\bibitem{fornengo2011}
N.~Fornengo, R.~Lineros, M.~Regis, and M.~Taoso,
``Possibility of a Dark Matter Interpretation for the Excess in Isotropic Radio
Emission Reported by ARCADE,''
\textit{Phys.\ Rev.\ Lett.} \textbf{107}, 271302 (2011).

\bibitem{fang2016}
K.~Fang and T.~Linden,
``Cluster Mergers and the Origin of the ARCADE-2 Excess,''
\textit{J.\ Cosmol.\ Astropart.\ Phys.} \textbf{2016}(10), 004 (2016).

\bibitem{digennaro2021}
G.~Di~Gennaro, R.~J. van~Weeren, F.~de~Gasperin, et al.,
``Fast magnetic field amplification in distant galaxy clusters,''
\textit{Nature Astron.} \textbf{5}, 268--275 (2021).

\bibitem{botteon2022}
A.~Botteon, T.~W. Shimwell, R.~Cassano, et al.,
``Magnetic fields and relativistic electrons fill entire galaxy clusters,''
\textit{Sci.\ Adv.} \textbf{8}, eabq7623 (2022).

\bibitem{botteon2020bridge}
A.~Botteon, R.~J. van~Weeren, G.~Brunetti, et al.,
``A giant radio bridge connecting two galaxy clusters in Abell 1758,''
\textit{Mon.\ Not.\ R.\ Astron.\ Soc.\ Lett.} \textbf{499}, L11--L15 (2020).

\bibitem{govoni2019}
F.~Govoni, E.~Orrù, A.~Bonafede, et al.,
``A radio ridge connecting two galaxy clusters in a filament of the cosmic web,''
\textit{Science} \textbf{364}, 981--984 (2019).

\bibitem{giovannini2010}
G.~Giovannini, A.~Bonafede, L.~Feretti, F.~Govoni, and M.~Murgia,
``The diffuse radio filament in the merging system ZwCl 2341.1+0000,''
\textit{Astron.\ Astrophys.} \textbf{511}, L5 (2010).

\bibitem{widrow2002}
L.~M. Widrow,
``Origin of galactic and extragalactic magnetic fields,''
\textit{Rev.\ Mod.\ Phys.} \textbf{74}, 775--823 (2002).

\bibitem{durrer2013}
R.~Durrer and A.~Neronov,
``Cosmological magnetic fields: their generation, evolution and observation,''
\textit{Astron.\ Astrophys.\ Rev.} \textbf{21}, 62 (2013).

\bibitem{jedamzik2025}
K.~Jedamzik, T.~Vachaspati, et al.,
``Hints of primordial magnetic fields at recombination and implications for the
Hubble tension,''
\textit{Nature Astron.} (2025).

\bibitem{williams1971}
E.~R. Williams, J.~E. Faller, and H.~A. Hill,
``New Experimental Test of Coulomb's Law: A Laboratory Search for Long-Range
Deviations,''
\textit{Phys.\ Rev.\ Lett.} \textbf{26}, 721--724 (1971).

\end{thebibliography}

\end{document}
